\documentclass[11pt]{article}
\usepackage[margin=1in]{geometry}
\usepackage{amsmath,amssymb,amsthm,booktabs,graphicx,hyperref}
% \usepackage{lineno}

\newtheorem{theorem}{Theorem}[section]
\newtheorem{lemma}[theorem]{Lemma}
\newtheorem{conjecture}[theorem]{Conjecture}
\newtheorem{remark}[theorem]{Remark}

\title{Spectral Classes of Polynomial Continued Fractions:\\A Hypergeometric Framework and Arithmetic Obstructions}
\author{Papanokechi\\
\small Independent Researcher, Yokohama, Japan\\
\small ORCID: 0009-0000-6192-8273}
\date{April 2026}

\begin{document}
\maketitle
% \linenumbers

\begin{abstract}
We propose a spectral classification program for polynomial continued fractions (PCFs) in which each family is assigned the invariant
\[
\operatorname{Spec}(K)=(d,\Lambda,\Delta,\rho,\tau),
\]
recording the dominant degree, normalized Poincar\'e root, discriminant class, Wallis tail-decay rate, and arithmetic sector. This framework organizes classical logarithmic and $1/\pi$ ladders together with higher-rank obstructed examples inside a single reproducible taxonomy. Our first proved result is the Poincar\'e characteristic lemma, showing that balanced degree-one PCFs are governed at leading order by the quadratic law $\Lambda^2-\alpha\Lambda-\delta=0$. Our second proved result is the Catalan logarithmic Frobenius obstruction lemma: after factorial normalization, a universal residual term $3/(2n)$ forces a logarithmic correction and rules out the naive Gauss ${}_2F_1$ ansatz. We then isolate two main conjectures, one for a clean $d=2$ V\_quad regime linked to a confluent-Heun/Stokes problem and one for an Ap\'ery-type $d=3$ regime governed by a Picard--Fuchs operator and numerically identified with $6/\zeta(3)$. The current dataset contains 37 rigorously verified families, including seven rows imported from Cohen (2024) that supply elliptic and Gauss-hypergeometric comparison examples. The resulting picture suggests that arithmetic behavior in PCFs is controlled not only by coarse spectral data, but also by Frobenius resonance, monodromy type, and higher-rank geometry. Key open problems are to prove nonreducibility in the clean obstructed sectors and to locate further resonant-versus-clean families.
\end{abstract}

\noindent\textbf{Keywords:} polynomial continued fractions, spectral classification, hypergeometric functions, arithmetic obstructions, Poincar\'e characteristic roots

\noindent\textbf{2020 Mathematics Subject Classification.}
Primary: 11J70;
Secondary: 33C05, 11J72, 34M35

\section{Introduction}
Polynomial continued fractions (PCFs) sit at a crossroads of special functions, Diophantine approximation, and arithmetic geometry \cite{wall1948,lorentzen2008}. Classical continued fractions of Gauss, Euler, Brouncker, and Ap\'ery show that low-degree recurrences can encode logarithms, Gamma quotients, elliptic periods, and zeta-type constants with remarkable efficiency \cite{apery1979}. Recent large-scale search programs, especially the Ramanujan Machine paradigm, have greatly expanded the list of examples \cite{raayoni2021}, but they also expose a conceptual gap: there is still no generally accepted invariant that predicts when a polynomial continued fraction should belong to a classical hypergeometric sector and when it should instead fall into a genuinely obstructed arithmetic class.

The central claim of this paper is that the first layer of that answer is spectral. Given a PCF with polynomial data $a_n,b_n$, we attach the invariant
\[
\operatorname{Spec}(K)=(d,\Lambda,\Delta,\rho,\tau),
\]
where $d$ records dominant growth, $\Lambda$ is the normalized Perron--Poincar\'e root, $\Delta$ is the balanced discriminant, $\rho$ is a Wallis-style tail-decay rate, and $\tau$ is the observed arithmetic sector. The philosophy is to separate coarse asymptotic behavior from fine arithmetic behavior. The pair $(d,\Lambda,\Delta)$ organizes families into broad sheets; the additional parameters $\rho$ and $\tau$, together with the fitted hypergeometric point $z_0$ and the Frobenius type of the normalized recurrence, refine those sheets into a workable taxonomy.\footnote{A companion paper \cite{pcf_bifurcation_2026} establishes the bifurcation phenomenon for specific families; the present work provides the general framework.} The broader synthesis also interfaces with our earlier PCF and asymptotic manuscripts \cite{pcf_unified_2026,ratio_universality_2026}.

\begin{remark}[Series position]
This paper is the third in a series. Paper~\cite{ratio_universality_2026} establishes ratio universality for Meinardus-class partition functions. Paper~\cite{pcf_bifurcation_2026} identifies the bifurcation phenomenon for specific PCF sectors. The present work provides the unifying spectral classification framework. No theorem in this paper duplicates a result in~\cite{ratio_universality_2026,pcf_bifurcation_2026}.
\end{remark}

This viewpoint is motivated by three distinct strands of evidence. First, the classical degree-one ladders for logarithms and rational multiples of $1/\pi$ share closely related recurrences yet land in different arithmetic sectors. Second, the Catalan family shows that matching the leading Poincar\'e root is not enough: a repeated-root Frobenius defect can force logarithmic behavior and destroy the naive Gauss ${}_2F_1$ ansatz, in the sense of the Frobenius local theory \cite{frobenius1873}. Third, the V\_quad and Ap\'ery rows display clean higher-rank asymptotics with no logarithmic resonance, but still resist collapse to the low-rank hypergeometric templates. Put differently, spectral data explain \emph{where} a family lives, while Frobenius resonance, monodromy, and Picard--Fuchs/Heun structure explain \emph{why} it does or does not reduce to a familiar closed form.

Our first rigorous result is a leading-order classification theorem in the balanced degree-one regime, in the spirit of the asymptotic recurrence theory initiated by Poincar\'e \cite{poincare1885}. It shows that the normalized ratio $r_n/n$ cannot converge arbitrarily: its limit is forced by the top coefficients of the recurrence. This theorem supplies the basic invariant used throughout the Phase~1 table and anchors the entire numerical taxonomy. Our second rigorous result is a concrete obstruction theorem for Catalan's constant. After factorial normalization, the symbolic residual contains an unavoidable term $3/(2n)$; since this coefficient cannot be removed inside a pure Frobenius power-law ansatz, a logarithmic correction is forced. This produces a clean proof-of-concept that arithmetic obstruction can be detected directly from normalized asymptotics.

Two further sectors emerge clearly from the data but remain conjectural. The first is the V\_quad regime, a clean $d=2$ family whose renormalized characteristic polynomial has distinct roots and whose Borel transform points to a confluent-Heun / Stokes connection problem rather than a Gauss-type evaluation; see \cite{ronveaux1995}. The second is the Ap\'ery regime, a balanced $d=3$ family governed by the classical Picard--Fuchs operator for the Ap\'ery recurrence and numerically identified with $6/\zeta(3)$; see Section~\ref{subsec:apery-outline}. These examples suggest a three-way obstruction trichotomy: repeated-root logarithmic resonance (Catalan), clean $d=2$ Heun/Stokes behavior (V\_quad), and clean $d=3$ Picard--Fuchs/modular behavior (Ap\'ery).

For ease of reference, we record the four principal statements in the front matter.

\begin{theorem}[Poincar\'e characteristic law, informal]
In the balanced degree-one class, the normalized root $\Lambda$ is forced by the quadratic equation $\Lambda^2-\alpha\Lambda-\delta=0$; see Theorem~\ref{thm:poincare-outline}.
\end{theorem}

\begin{lemma}[Catalan Frobenius obstruction, informal]
For the Catalan row, factorial normalization leaves a residual $3/(2n)$ term, forcing a logarithmic correction and excluding the plain Gauss ${}_2F_1$ model; see Lemma~\ref{lem:catalan-outline}.
\end{lemma}

\begin{conjecture}[V\_quad clean $d=2$ obstruction, informal]
The V\_quad constant belongs to a confluent-Heun / Stokes sector and is not reducible to a single algebraic Gauss value; see Conjecture~\ref{conj:vquad-outline}.
\end{conjecture}

\begin{conjecture}[Ap\'ery clean $d=3$ obstruction, informal]
The Ap\'ery row belongs to a Picard--Fuchs / modular sector, distinct both from the Catalan resonance mechanism and from the V\_quad Heun regime; see Conjecture~\ref{conj:apery-obstruction} and Section~\ref{subsec:apery-outline}.
\end{conjecture}

The current taxonomy contains \textbf{37 rigorously verified families}. It includes ten $d=2$ rows, five elliptic examples, and a seven-row external block mined from Cohen's 2024 database of polynomial-type continued fractions \cite{cohen2024database}, whose data supply especially useful comparison cases in the elliptic and Gauss-hypergeometric sectors. The visual appendix makes this classification immediately visible: the classical logarithmic and $1/\pi$ rows occupy the low-degree mainland, the Cohen additions enrich the elliptic comparison layer, and the V\_quad family remains isolated by its unusually large convergence-rate parameter $\rho$; see Appendix~\ref{app:visuals}. These plots complement the table by showing not just exact values but also convergence profiles and sector separation.

The structure of the paper follows the discovery pipeline. Section~2 fixes definitions and normalization conventions. Section~3 proves the Poincar\'e characteristic lemma. Section~4 presents the 37-family taxonomy table and records the Cohen additions. Section~5 develops the $z_0$-orbit framework for the degree-one hypergeometric sector. Section~6 studies the obstruction mechanisms, moving from the proved Catalan resonance lemma to the conjectural clean sectors for V\_quad and Ap\'ery. Section~7 collects open problems, including nonreducibility in the clean sectors and the search for further resonant-versus-clean examples. The appendix records the scripts, figures, and data products that make the workflow reproducible.

\section{Definitions and spectral classes}
\subsection{Polynomial continued fractions}
Define
\[
K(a_n,b_n) = b_0 + \cfrac{a_1}{b_1 + \cfrac{a_2}{b_2 + \cdots}}
\]
with polynomial data $a_n,b_n$.

\subsection{The spectral invariant}
Introduce
\[
\operatorname{Spec}(K)=(d,\Lambda,\Delta,\rho,\tau),
\]
where $d$ is the growth degree, $\Lambda$ is the Poincar\'e root, $\Delta$ the discriminant, $\rho$ the Wallis decay rate, and $\tau$ the arithmetic type.

\begin{remark}
All numerical invariants quoted in this outline are extracted from the reproducible CSV/JSON outputs generated by the workspace scripts. In particular, $\Lambda$ is estimated from normalized convergent ratios, $\rho$ from the Wallis-determinant tail slope, and the arithmetic tag $\tau$ from high-precision PSLQ and special-function matching.
\end{remark}

\section{The Poincar\'e lemma}
\begin{theorem}[Poincar\'e characteristic lemma]\label{thm:poincare-outline}
Let $p_n$ satisfy
\[
  p_n=b_n p_{n-1}+a_n p_{n-2},
\]
with
\[
  a_n = \delta n^2 + O(n), \qquad b_n = \alpha n + O(1),
\]
and suppose the normalized ratio $r_n:=p_n/p_{n-1}$ obeys $r_n/n\to\Lambda\neq 0$. Then $\Lambda$ satisfies the characteristic equation
\[
  \Lambda^2 - \alpha \Lambda - \delta = 0.
\]
In particular, the balanced degree-one spectral class is already determined at leading order by the pair $(\alpha,\delta)$.
\end{theorem}

\begin{proof}[Proof sketch]
From the recurrence one has the exact identity
\[
  r_n = b_n + \frac{a_n}{r_{n-1}}.
\]
Divide by $n$ and use the hypotheses $b_n/n \to \alpha$, $a_n/n^2 \to \delta$, and $r_{n-1}/n \to \Lambda$. Passing to the limit gives
\[
  \Lambda = \alpha + \frac{\delta}{\Lambda},
\]
which is equivalent to the quadratic equation above. In the computational pipeline this lemma is used in reverse: one first estimates $\Lambda$ numerically from convergent ratios and then checks the resulting discriminant class against the symbolic leading coefficients.
\end{proof}

\section{Phase 1: Taxonomy table}
\subsection{Seed families}
Include at least 15 families, for example:
\begin{itemize}
  \item logarithmic ladder `log-k`,
  \item pi-family `pi-m`,
  \item Catalan-type degenerate cases,
  \item V\_quad and Ap\'ery-type outliers,
  \item probe families used to test boundary and sign hypotheses.
\end{itemize}

\subsection{Main taxonomy table}
The current CSV now contains \textbf{37 verified rows}. A representative submission-ready slice is shown below, with the seven Cohen-derived additions highlighted alongside the earlier seed set.

\begin{center}
\scriptsize
\begin{tabular}{@{}lcccccl@{}}
\toprule
Family & $d$ & $\Lambda$ & $\Delta$ & $\rho$ & $\tau$ & Status \\
\midrule
log-k=2 & 1 & 2 & 1 & 0.301 & log-type & seed \\
log-k=3 & 1 & 3 & 4 & 0.477 & log-type & seed \\
log-k=4 & 1 & 4 & 9 & 0.602 & log-type & seed \\
log-k=5 & 1 & 5 & 16 & 0.699 & log-type & seed \\
log-k=6 & 1 & 6 & 25 & 0.779 & log-type & seed \\
log-k=7 & 1 & 7 & 36 & 0.845 & log-type & seed \\
log-k=8 & 1 & 8 & 49 & 0.903 & log-type & seed \\
\textbf{log-k=9} & 1 & 9 & 64 & 0.955 & log-type & \textbf{NEW} \\
log-k=10 & 1 & 10 & 81 & 1.000 & log-type & seed \\
\textbf{log-k=11} & 1 & 11 & 100 & 1.042 & log-type & \textbf{NEW} \\
\textbf{log-k=12} & 1 & 12 & 121 & 1.080 & log-type & \textbf{NEW} \\
pi-m=0 & 1 & 2 & 1 & 0.301 & pi-type & seed \\
pi-m=1 & 1 & 2 & 1 & 0.302 & pi-type & seed \\
pi-m=3 & 1 & 2 & 1 & 0.305 & pi-type & seed \\
\textbf{pi-m=6} & 1 & 2 & 1 & 0.308 & pi-type & \textbf{NEW} \\
\textbf{pi-m=7} & 1 & 2 & 1 & 0.309 & pi-type & \textbf{NEW} \\
\textbf{cohen-4D1-k=1/2} & 2 & 4 & 9 & 0.602 & elliptic-period & \textbf{COHEN} \\
\textbf{cohen-4D1.7-k=1/2} & 2 & 4 & 9 & 0.602 & elliptic-period & \textbf{COHEN} \\
\textbf{cohen-4D3.3-k=1/2} & 3 & 2 & $1/4$ & 0.125 & elliptic-period & \textbf{COHEN} \\
\textbf{cohen-4D4-k=1/2} & 1 & $3/2$ & 4 & 0.477 & elliptic-period & \textbf{COHEN} \\
\textbf{cohen-4D6-k=1/2} & 1 & 2 & $9/4$ & 0.602 & elliptic-period & \textbf{COHEN} \\
\textbf{cohen-5.2.7.4} & 1 & $1.707\ldots$ & 2 & 0.766 & ${}_2F_1$-type & \textbf{COHEN} \\
\textbf{cohen-3.1.11.3-z=1} & 2 & $11.657\ldots$ & 128 & 1.531 & log-type & \textbf{COHEN} \\
apery-zeta3 & 3 & $17+12\sqrt2$ & 1152 & 3.062 & Apery-isolated & seed \\
vquad-3-1-1 & 2 & 3 & -11 & 13.213 & V\_quad-mystery & obstruction \\
\bottomrule
\end{tabular}
\end{center}

\section{Phase 2: The z0-orbit framework}
\subsection{Hypergeometric fitting}
Explain the fitting ansatz
\[
L = C\,{}_2F_1(a,b;c;z_0)
\]
and the recovery of $(a,b,c)$ from $(\alpha,\delta)$.

\begin{theorem}[z0-orbit separation, empirical]\label{thm:z0-orbit-outline}
For the identified $d=1$ families in the current dataset, the $PGL_2(\mathbb{Q})$-orbit of the fitted evaluation point $z_0$ is a finer invariant than $(\Lambda,\Delta)$ and separates the mixed control pair \texttt{log-k=2} and \texttt{pi-m=0}.
\end{theorem}

\begin{proof}[Proof sketch]
The Phase~2 audit fits all currently identified $d=1$ rows to values of the form $L=C\,{}_2F_1(a,b;c;z_0)$ with small rational $C$. The control pair \texttt{log-k=2} and \texttt{pi-m=0} has the same indicial triple $(a,b,c)$ but different fitted values of $z_0$, and these points are not related by the Euler--Pfaff sixfold action. Thus $(a,b,c)$ alone is insufficient, while adjoining the orbit of $z_0$ restores separation.
\end{proof}

\section{Phase 3: Arithmetic obstructions}
\subsection{Catalan obstruction}
\begin{lemma}[Catalan obstruction]\label{lem:catalan-outline}
For the Catalan family
\[
  a_n=-4n^2,\qquad b_n=4n+1,
\]
one has $\Delta=0$ and repeated Poincar\'e root $\Lambda=2$. After normalizing by $u_n:=p_n/n!$, the recurrence becomes
\[
  u_n=(4+1/n)u_{n-1}-(4-4/n)u_{n-2}.
\]
No solution of the form $u_n=2^n n^s(1+c_1/n+c_2/n^2+\cdots)$ can satisfy this recurrence without a logarithmic correction term.
\end{lemma}

\begin{proof}[Proof sketch]
Substituting the nonlogarithmic Frobenius ansatz into the normalized recurrence gives the symbolic residual
\[
  \mathrm{RHS}-\mathrm{LHS}=\frac{3}{2n}+\frac{3c_1-2s^2-3s}{2n^2}+O(n^{-3}).
\]
The coefficient $3/(2n)$ is independent of $s,c_1,c_2,\dots$, so it cannot be removed inside the pure power-law ansatz. Hence the second local solution is necessarily logarithmic. This is the mechanism behind the observed failure of the plain Gauss ${}_2F_1$ fit for Catalan. A 220-digit PSLQ check against the standard algebraic/logarithmic basis also returns \texttt{None}, supporting the obstruction numerically.
\end{proof}

\subsection{V\_quad obstruction (conditional / numerical)}
\begin{lemma}[Clean d=2 tier for V\_quad]\label{lem:vquad-outline}
For the recurrence
\[
p_n = (3n^2+n+1)p_{n-1} + p_{n-2},
\]
one has the asymptotic law
\[
p_n \sim C\,3^n (n!)^2 n^{1/3},
\]
so the renormalized recurrence has distinct characteristic roots and no logarithmic Frobenius resonance.
\end{lemma}

\begin{proof}[Proof sketch]
Set
\[
  p_n = 3^n (n!)^2 n^{1/3} q_n.
\]
Then the recurrence becomes exactly
\[
  q_n = A_n q_{n-1} + B_n q_{n-2},
\]
where
\[
  A_n = \frac{(n-1)^{1/3}(3n^2+n+1)}{3n^{7/3}},
  \qquad
  B_n = \frac{(n-2)^{1/3}}{9n^{7/3}(n-1)^2}.
\]
Expanding at infinity gives
\[
  A_n = 1 + \frac{1}{9n^2} - \frac{17}{81n^3} - \frac{8}{81n^4} + O(n^{-5}),
  \qquad
  B_n = \frac{1}{9n^4} + \frac{4}{27n^5} + \frac{11}{81n^6} + O(n^{-7}).
\]
Hence the asymptotic characteristic equation is
\[
  \lambda^2 - \lambda = 0,
\]
with distinct roots $\lambda_1=1$ and $\lambda_2=0$. Therefore the V\_quad sector is a clean $d=2$ regime with one dominant power-law branch and one recessive branch, but no logarithmic Frobenius resonance.
\end{proof}

This strengthens the obstruction picture considerably. Catalan fails because of repeated-root Frobenius resonance; V\_quad fails for a different reason. After the factorial-square normalization the dominant solution stabilizes, while the second solution is recessive and numerically tiny, so generic initial data collapse onto the same dominant branch without any sign of a logarithmic defect.

\paragraph{Generating-function / Heun link.}
The naive ordinary generating function for the normalized sequence $q_n$ is not especially transparent because of the fractional factor $n^{1/3}$. However, the associated Birkhoff transform leads to an explicit second-order differential equation of confluent-Heun type. In the notation of the April 9 connection-matrix computation, the transformed function $y(x)$ satisfies
\[
  (3x^2+x+1)y'' + (6x+1)y' - x^2 y = 0.
\]
Its finite singularities are
\[
  x = \frac{-1 \pm i\sqrt{11}}{6},
\]
while $x=\infty$ is an irregular singularity of rank $1$. Thus the equation is not ordinary HeunG; rather, it lies in a confluent-Heun / Stokes sector, exactly matching the arithmetic behavior seen numerically.

Equivalently, if $F(z)=\sum_{n\ge 0} p_n z^n$ is the ordinary generating function for the raw denominator sequence, then direct summation of the recurrence yields the differential equation
\[
  3z^3 F''(z) + 10z^2 F'(z) + (z^2+5z-1)F(z) + 1 = 0,
\]
whose Birkhoff-Laplace transform is the equation above. This is the clearest analytic route currently available toward a full V\_quad theorem.

\paragraph{Connection to known special functions.}
The finite-singularity set has discriminant $-11$, which points naturally toward CM-type and modular/elliptic structures; nevertheless the high-precision tests performed so far do not collapse the value to a simple Gauss, elliptic, Bessel, Airy, or ${}_0F_2$ constant. The most plausible current interpretation is therefore: V\_quad is not hypergeometric in the plain ${}_2F_1$ sense, but is governed by a confluent-Heun-type connection problem with nontrivial Stokes data.

At present the evidence is numerical but strong: the 220-digit, 320-digit, and now 400-digit PSLQ searches against the expanded basis
\[
\begin{aligned}
\{&1, V_{\mathrm{quad}}, \pi, \log 2, \log 3, \sqrt 3, \sqrt{11},\\
&\Gamma(1/3), \Gamma(2/3), \Gamma(1/4), \Gamma(3/4),\\
&K(1/\sqrt 2), K(\sqrt 3/2), E(1/\sqrt 2), E(\sqrt 3/2), \pi^2, \zeta(3)\}
\end{aligned}
\]
all return \texttt{None}. Thus V\_quad is best stated as a conjectural arithmetic obstruction supported by reproducible Heun-linked evidence.

\begin{conjecture}[V\_quad arithmetic obstruction]\label{conj:vquad-outline}
The V\_quad constant is not representable by a single algebraic Gauss value $C\,{}_2F_1(a,b;c;z_0)$ and is more naturally associated with a confluent-Heun / Stokes-type sector.
\end{conjecture}

\subsection{Ap\'ery d=3 obstruction (conditional / modular)}\label{subsec:apery-outline}
The isolated Ap\'ery row from the taxonomy is governed by the recurrence
\[
  p_n = (34n^3+51n^2+27n+5)p_{n-1} - n^6 p_{n-2},
\]
with the experimentally verified value
\[
  V_{\mathrm{Ap}} = \frac{6}{\zeta(3)}.
\]
Let
\[
  \lambda_+ = 17+12\sqrt2, \qquad \lambda_- = 17-12\sqrt2 = \lambda_+^{-1},
\]
and normalize by
\[
  p_n = \lambda_+^n (n!)^3 n^{3/2} q_n.
\]
Then one obtains the exact recurrence
\[
  q_n = A_n q_{n-1} + B_n q_{n-2},
\]
where
\[
  A_n = \frac{(n-1)^{3/2}(34n^3+51n^2+27n+5)}{\lambda_+ n^{9/2}},
  \qquad
  B_n = -\frac{n^{3/2}(n-2)^{3/2}}{\lambda_+^2 (n-1)^3}.
\]
Expanding at infinity gives
\[
  A_n = \bigl(1+\lambda_+^{-2}\bigr) - \frac{147}{4\lambda_+ n^2} - \frac{57}{4\lambda_+ n^3} + O(n^{-4}),
\]
\[
  B_n = -\lambda_+^{-2} + \frac{3}{2\lambda_+^2 n^2} + \frac{3}{\lambda_+^2 n^3} + O(n^{-4}).
\]
Hence the limiting characteristic equation is
\[
  r^2 - \bigl(1+\lambda_+^{-2}\bigr)r + \lambda_+^{-2} = (r-1)(r-\lambda_+^{-2}) = 0,
\]
with distinct roots $1$ and $\lambda_+^{-2} = \lambda_-^2$. Thus the Ap\'ery sector is \emph{not} Catalan-like: there is no repeated-root Frobenius resonance and no forced logarithmic term at the dominant asymptotic level. In this sense it is a clean higher-rank obstruction, analogous to V\_quad but now in the balanced $d=3$ regime.

\paragraph{Picard--Fuchs form.}
After the factorial-cube Borel normalization, the companion generating function $U(z)$ is governed by the classical Ap\'ery operator
\[
  \Bigl[\theta^3 - z(2\theta+1)(17\theta^2+17\theta+5) + z^2(\theta+1)^3\Bigr]U(z)=0,
  \qquad \theta = z\frac{d}{dz}.
\]
This is a third-order Fuchsian Picard--Fuchs equation, not a plain Gauss ${}_2F_1$ equation and not an ordinary HeunG equation. The natural arithmetic interpretation is therefore modular / Ap\'ery-type rather than degree-one hypergeometric or degree-two Heun-like.

\paragraph{420-digit PSLQ evidence.}
At 420 digits, PSLQ on the expanded basis
\[
\begin{aligned}
\{&1, V_{\mathrm{Ap}}, \pi, \log 2, \log 3, G, \sqrt2, \sqrt3, \sqrt5, \sqrt{11},\\
&\Gamma(1/3), \Gamma(2/3), \Gamma(1/4), \Gamma(3/4), K(1/\sqrt2), E(1/\sqrt2), \pi^2, \pi^3\}
\end{aligned}
\]
returns \texttt{None}; only after adjoining $1/\zeta(3)$ does PSLQ recover the exact relation
\[
  -V_{\mathrm{Ap}} + \frac{6}{\zeta(3)} = 0.
\]
So the obstruction is not one of simple nonexistence of a closed form, but rather that the correct closed form lies in a genuinely Ap\'ery / Picard--Fuchs sector outside the lower-rank Gauss and Heun templates.

\begin{conjecture}[Ap\'ery modular obstruction]\label{conj:apery-obstruction}
The Ap\'ery row is not reducible to a single algebraic Gauss value $C\,{}_2F_1(a,b;c;z_0)$, nor to the clean $d=2$ confluent-Heun sector of V\_quad. Its correct home is a balanced $d=3$ Picard--Fuchs / modular class.
\end{conjecture}

\section{Conjectures and open problems}
\begin{itemize}
  \item Extend the taxonomy beyond the current seed set.
  \item Determine whether every arithmetic sector is controlled by a hypergeometric orbit, a logarithmic Frobenius defect, or a higher-rank Stokes datum.
  \item Prove or disprove the V\_quad obstruction conjecture.
  \item Develop a full obstruction theory for $d \ge 2$ families.
\end{itemize}

\section{Appendix: verification scripts}\label{app:verification}
\begin{itemize}
  \item \texttt{spectral\_class\_phase1.py}
  \item \texttt{spectral\_class\_phase2\_congruence.py}
  \item \texttt{vquad\_connection\_matrix.py}
  \item \texttt{vquad\_heun\_connection\_check.py}
  \item \texttt{results/generate\_spectral\_visualizations.py}
\end{itemize}

\subsection{Visualizations of Spectral Classes}\label{app:visuals}
The appendix figures provide a compact visual summary of the 37-family taxonomy and make the two main structural messages easy to see at a glance: the Cohen elliptic block forms a coherent low-$d$ cluster, while the V\_quad rows remain isolated by their unusually large convergence-rate parameter $\rho$.

\begin{figure}[htbp]
\centering
\includegraphics[width=0.88\linewidth]{figures/taxonomy_rho_vs_d.png}
\caption{Scatter plot of the estimated Wallis convergence rate $\rho$ against spectral degree $d$, colored by arithmetic sector. The highlighted labels mark the V\_quad obstruction candidate and the newest Cohen additions, including the elliptic specializations \texttt{cohen-4D1-k=1/2}, \texttt{cohen-4D1.7-k=1/2}, \texttt{cohen-4D4-k=1/2}, \texttt{cohen-4D6-k=1/2}, and the direct Gauss row \texttt{cohen-5.2.7.4-c=3/2-z=1/2}.}
\label{fig:taxonomy-rho-vs-d}
\end{figure}

\begin{figure}[htbp]
\centering
\includegraphics[width=0.84\linewidth]{figures/taxonomy_lambda_distribution.png}
\caption{Distribution of the normalized root $\Lambda$ across the main regimes. The $d=1$ balanced sector contains the logarithmic and $1/\pi$ ladders; the $d=2$ sector contains both the V\_quad obstruction family and the new Cohen logarithmic/elliptic witnesses; the $d=3$ sector includes the Ap\'ery and elliptic-square examples.}
\label{fig:taxonomy-lambda-distribution}
\end{figure}

\begin{figure}[htbp]
\centering
\includegraphics[width=0.88\linewidth]{figures/convergence_representatives.png}
\caption{Representative convergence curves $\log_{10}(\mathrm{error})$ versus depth for a log-ladder row, a $\pi$-family row, the V\_quad obstruction candidate, one Cohen elliptic example, and one Cohen ${}_2F_1$ specialization. The contrast emphasizes both the stability of the classical hypergeometric sectors and the numerically distinctive behavior of V\_quad.}
\label{fig:convergence-representatives}
\end{figure}

\section*{Data Availability}
All numerical data supporting the results of this study were generated 
computationally using Python/mpmath and are fully reproducible from the 
algorithms described in the manuscript. Computation scripts are publicly 
available at \url{https://github.com/papanokechi/pcf-spectral-classes}.

\section*{Conflict of Interest}
The author declares no conflicts of interest.

\section*{Funding}
This research received no external funding. The work was conducted 
independently by the author.

\section*{Acknowledgments}
We acknowledge Henri Cohen for making the 2024 polynomial-continued-fraction database publicly available; the Cohen block now supplies a valuable elliptic and Gauss-hypergeometric comparison layer for the present taxonomy.

Computational work in this paper was conducted using the SIARC
v6.5 pipeline with AI assistance from GitHub Copilot and
Anthropic Claude. All mathematical claims and conclusions are
the author's own.

\bibliographystyle{amsplain}
\bibliography{pcf_spectral_refs}

\end{document}
